\section{Methodology}
\label{sec:methodology}

Half of what this project has produced is the apparatus, and it is presented here as method rather
than as housekeeping, because every component of it was built in response to a specific measured
failure --- most of them this project's own. A reader working in another game can adopt these
pieces individually; each is stated with the failure it prevents.

\subsection{The duplicate-block design}
\label{sec:methodology-duplicate}

\textbf{What it is.} A \emph{cell} is a comparison of two configurations $A$ and $B$ over a
\emph{bank} of deals identified by one integer seed. Every deal in the bank is played twice: once
with $A$ in the odd seats and $B$ in the even seats, and once with the seats swapped. Both
configurations therefore play the identical deal from both sides.

\textbf{What it prevents.} Fish has enormous per-deal variance --- a deal can be won or lost on the
cards --- and an unpaired comparison spends most of its sample resolving that variance rather than
the difference between the policies. The rotation cancels the deal's intrinsic luck exactly, in the
same spirit as duplicate bridge and the poker line's variance-reduction estimators
\cite{zinkevich-divat,burch-aivat}.

\textbf{What it requires.} Reproducibility from the seed alone (\S\ref{sec:engine-impl}). If the
engine's randomness is not keyed by \emph{(seed, seat, purpose)}, the two rotations are not the same
deal and the design silently degrades into an unpaired one.

Throughout this paper a cell's result is reported as an \textbf{edge} in points:
$100 \times (\text{winRate}_A - 0.5)$, where $0$ means the two configurations are indistinguishable.

\subsection{Deal-clustered intervals}
\label{sec:methodology-intervals}

The two rotations of one deal are \emph{not} two independent observations. They share a deal, and
their outcomes are strongly negatively correlated by construction. Treating them as independent ---
which is what a Wilson interval on the game count does --- understates the width.

Every interval in this paper is a \textbf{deal-clustered bootstrap}: 20{,}000 resamples in which the
resampling unit is the \emph{deal}, carrying both of its rotations together. No Wilson interval is
quoted anywhere in this paper, and the substitution is not cosmetic: it is the difference between a
claim that replicates and one that does not.

\paragraph{Pooling two banks.} Two banks of equal size pool as the mean of the two edges with
$\text{se}_{\text{pooled}} = \sqrt{\text{se}_1^2 + \text{se}_2^2}/2$, where each $\text{se}$ is read
back off that bank's clustered interval. Both per-bank values are printed beside every pooled figure
in this paper, so a reader can always see the pooling rather than take it on trust.

\paragraph{Differences of cells.} A difference between two \emph{pooled} cells --- for example an
ablation's leave-one-out drop --- is the difference of the estimates with the two pooled half-widths
combined \textbf{in quadrature}. This is conservative. The arms are paired on deals, so a genuine
paired delta would be narrower, but the harness produces no paired delta \emph{across} two separately
run cells, and inventing one would be a stronger claim than the data structure supports.

\subsection{Power, computed in advance}
\label{sec:methodology-power}

At parity a cell's 95\% half-width in points is approximately $98/\sqrt{N}$ for $N$ games. The
harness computes this for every cell and emits it with the result, including the mirror case where
no half-width exists.

The arithmetic is stated in advance because it decides what may be attempted. Resolving one point
needs roughly \vsevenScoreboardGames{} games; a quarter of a point needs sixteen times that. A study
that wants to find quarter-point mechanisms on the scoreboard cannot afford to, and
\S\ref{sec:towardone-optimal} names that as an open problem rather than pretending it is solved.

\textbf{The failure this prevents} is the one \S\ref{sec:arc-vsix-apparatus} describes: five
mechanisms that cleared a 95\% paired interval at one budget and returned exactly zero at three
times it. Sample sizes in this paper are chosen from the arithmetic before the cell is run, never
after seeing it.

\subsection{Replication as a veto}
\label{sec:methodology-replication}

Every claim in this paper is measured on \emph{two disjoint banks}, and the rule is absolute:

\begin{quote}
A claim whose \textbf{sign} does not agree on both banks is reported as NOT REPLICATED, whatever its
pooled interval says.
\end{quote}

This is stricter than a pooled significance test and deliberately so. A pooled interval can exclude
zero while the two banks disagree, and in this corpus that configuration has occurred and has been
published as a positive result before. \S\ref{sec:against-location} is where the rule bites against
this paper's own release, and \S\ref{sec:results-panel} is where it bites against this paper's own
headline cell.

\subsection{Sealed material and preregistration}
\label{sec:methodology-seal}

\textbf{The problem.} A programme that generates candidate configurations, scores them, and picks
the best is performing a maximisation. The winner's score is upward-biased by exactly the amount a
maximum over $K$ draws exceeds the mean, and no amount of care in the final measurement removes it
if the final measurement uses material the selection has seen.

\textbf{The apparatus.} Before any candidate existed, deal banks were reserved by seed and their
digests committed: a rolling hash over the six dealt hands and the dealer of every deal in the bank,
computed without constructing a policy or playing a game. An adversary population was split in half
by a rule fixed in advance, and one half was \emph{sealed} --- encrypted, with the plaintext's
SHA-256 committed --- so that it could not be trained against even accidentally. The engine enforces
the seal: a run that names a sealed seed fails unless an explicit unseal variable is set, which is
set once, at the top of the final evaluation, and recorded.

\textbf{The protocol.} The entire final battery --- every cell, its sample size, its threshold, the
arithmetic behind the sample size, and the conditions that would count as failure --- was written
down and committed before any holdout bank had been played and before the sealed half had been
decoded. The evaluation phase then ran from a cleared context, reading that protocol and nothing
else from the project's document set. Where it needed a training figure for comparison, it took it
from the protocol, which states its expectations in advance for exactly that purpose, so that
agreement is a check and disagreement is a finding.

\textbf{Deviations.} Anything the evaluation did that the protocol does not specify is recorded as a
deviation, next to the affected result rather than in a footnote. There are \vsevenDeviations{}
(Appendix~\ref{app:deviations}). A battery with recorded deviations and honest numbers is worth more
than one that silently matched its own protocol.

\subsection{The commit gate runs before any strength number}
\label{sec:methodology-gate}

\textbf{The failure this prevents} is dissociation between win rate and soundness. This corpus
contains a configuration that scores \vsevenScoreboardTrapWin\% against a prior version while
carrying \vsevenScoreboardTrapDead\% provably-dead asks --- asks for cards its own team already
demonstrably owns --- a dead run of hundreds of consecutive such asks, and
\vsevenScoreboardTrapLimit\% of its games killed by the action cap. It is stronger and it is
broken.

So a configuration is gated \emph{before} any strength number is computed for it, on eight rules
covering dead asks, dead runs, action-limit games, the mirror game-length tail, late declarations,
and the side-channel tests below. A configuration that fails is rejected regardless of its win rate.
The gate is run on \emph{training} material by design: it is a soundness check on self-play
behaviour, not a strength claim.

\textbf{A gate needs a negative control.} The corpus maintains a configuration built to be
rejected, and the gate must reject it; a gate that cannot would certify nothing about the one it
accepts. \S\ref{sec:protocol-gate} reports both.

\subsection{Detection floors}
\label{sec:methodology-floor}

An exploitability number is produced by a search, and a search that finds nothing has two possible
explanations: there is nothing, or the search is too weak. A \textbf{detection floor} separates
them. Plant an artificial handicap of known size in a target, fit a responder against the
handicapped target, and measure the \emph{excess} of that responder's edge over its edge against the
unhandicapped one. The smallest planted edge the procedure recovers is the floor.

Two findings from doing this are methodologically important and neither was expected.

\begin{enumerate}
\item \textbf{The floor does not buy down with evaluation games.} An early prediction that it would
      fall as (evaluation games)$^{-1/2}$ was refuted at its own predicted value: at four times the
      games the floor is \vsevenPrimaryFloor{}, and the sub-floor rung is still undetected --- not
      because the interval is wide but because the excess resolves \emph{to} zero. Below roughly a
      point and a half the responder is \emph{empty}, not unresolved. The binding constraint is
      search power, not evaluation power.
\item \textbf{Floors are family-specific.} The declaration channel has its own, higher floor
      (\vsevenDeclFloor{}), and reading a declaration-objective result against the general floor
      would overstate it.
\end{enumerate}

A class of adversary may contribute to an exploitability claim only after it has passed a
planted-edge calibration at or below the effect size being claimed. That precondition is why the
negative controls of \S\ref{sec:results-controls} are run at all: an instrument that cannot recover
a planted edge of a given size cannot report the absence of a real one of that size.

\subsection{The homogeneity and side-channel gates}
\label{sec:methodology-side}

\textbf{The problem.} The team under evaluation is three identical agents with no shared secret.
``Identical and deterministic'' means their coordination is common knowledge
(\S\ref{sec:domain-team}) --- which is legitimate --- but it also means an implementation detail can
become an illegal signalling channel without anyone intending it. Before this cycle the constraint
had only ever been checked by reading the source.

\textbf{The apparatus.} Four mechanical tests, run over all four engine decision types:

\begin{description}
\item[S3, listening substitution.] Does a seat's behaviour change more when a \emph{teammate}'s
  hidden action changes than when an \emph{opponent}'s does, beyond what the rules explain?
\item[S4, stream independence.] Is the seat's private random stream drawn independently of the deal?
\item[S5, posterior invariance.] Does the seat's action depend on the true card location beyond what
  its own posterior supports?
\item[S6, seat isolation.] Replay every decision from (own hand, public event stream, rules, reset
  seed) alone and count the decisions the reconstruction cannot reproduce. Zero tolerance.
\end{description}

\textbf{The calibration, which is the part that makes it a test.} Three positive controls were
built, each a policy carrying one planted cheat, and each must fail exactly the tests the protocol
names for it in advance and no others. \S\ref{sec:protocol-material} reports the result on holdout
material. Without those controls, a clean report from the gate would be uninformative.

\textbf{One honest limitation, stated wherever the gate is used.} S6's zero-tolerance form requires
one thread --- the audit depends on the deal-to-thread partition, and work-stealing does not fix that
partition at any higher count --- and agents rebuilt per deal. Rebuilding agents per deal
\emph{changes play}. So the gate certifies ``with cross-deal residue removed, nothing else depends on
anything it should not''. The residue is not removed from the shipped mode, and every strength number
in this paper is measured under it (\S\ref{sec:results-residual}).

\subsection{Worst case and minimax regret over a shared panel}
\label{sec:methodology-panel}

The project's standing reporting rule is that \textbf{the worst case leads and an aggregate never
does}. A mean over an opponent panel rewards beating weak opponents by more, which is not the
quantity anyone cares about.

Two statistics are reported, both over a panel scored identically for every arm:

\begin{itemize}
\item the \textbf{worst cell} per arm --- the panel member on which it does least well;
\item the \textbf{minimax regret} per arm --- for each member, how far the arm falls short of the
      best arm on that member, maximised over members.
\end{itemize}

They can disagree, and in \S\ref{sec:results-panel} they do. Reporting both, and reporting the panel
mean last and labelled as a diagnostic, is the discipline; choosing whichever one flatters the
release is the failure it prevents. Every arm must carry a cell for every member, or minimax regret
is not defined at all --- which is why the panel is fixed by name in advance and scored for all arms
on the same banks.

\paragraph{A caveat the panel design forces.} An arm that is itself a panel member meets itself and
carries a free $0.00$ cell, which can only flatter its worst case. Where that matters,
\S\ref{sec:results-panel} prints the worst case both with and without self-cells.
